\documentclass[11pt,a4paper]{article}

\usepackage[margin=1in]{geometry}
\usepackage{amsmath}
\usepackage{amssymb}
\usepackage{booktabs}
\usepackage{longtable}
\usepackage{hyperref}

\hypersetup{colorlinks=true, linkcolor=black, urlcolor=blue, citecolor=black}

\newcommand{\cacose}{\texttt{CaCoSE}}
\newcommand{\caef}{\textit{CaEF}}

\title{\vspace{-2em}CaCoSE Phase 1: Benchmark Reproduction Record\\
\large Experimental log for the 27 August 2026 cluster run}
\author{GNN\_CaCoSE project}
\date{27 August 2026}

\begin{document}
\maketitle
\vspace{-1.5em}

\begin{abstract}
\noindent This is the laboratory record of the run that produced the Phase~1 reproduction
numbers: what was executed, on what hardware, under which commit and container, and what every
individual seed returned. It is deliberately separate from the design specification
(\texttt{phase1\_implementation.tex}, which states \emph{what} is built and why) and from
\texttt{REPRO\_REPORT.md} (which states \emph{what we conclude}). This document is the evidence
those two rest on, written so that a reader who has neither the cluster nor the raw logs can still
audit the claims, and so that the run can be repeated exactly. Where the record is incomplete or a
measurement is weaker than it appears, that is stated rather than omitted.
\end{abstract}

\section{Provenance}

Every number in this document comes from a single batch of 30 SLURM tasks executed on
27~August~2026. Nothing here is a local run, an average across runs, or a figure carried over from
an earlier attempt.

\begin{table}[!htbp]
\centering
\begin{tabular}{ll}
\toprule
Field & Value \\
\midrule
Commit & \texttt{2147b500cd12} \\
Cluster & AART \texttt{pleiades}, partition \texttt{pleiades} \\
Node & \texttt{pleiades-1-3} (all 30 tasks) \\
GPU & NVIDIA GeForce GTX 1080 Ti, one per task \\
Container & \texttt{/data/\$USER/CaCoSE/containers/cacose.sif} \\
Wall-clock window & 15:33:30 -- 15:45:18 EDT, 27 August 2026 \\
Python / torch & 3.11 / 2.5.1+cu121 \\
PyG / NumPy / SciPy & 2.6.1 / 1.26.4 / 1.17.1 \\
\bottomrule
\end{tabular}
\caption{Run provenance. The commit is recorded inside every result JSON by the runner, resolved on
the host before entering the container, since the container has no \texttt{git}.}
\end{table}

All 30 tasks landed on the same node. That is convenient for comparability -- no cross-node
hardware variation enters the seed spread -- but it was not enforced, and a rerun may distribute
differently.

\subsection{Jobs}

\begin{table}[!htbp]
\centering
\begin{tabular}{llll}
\toprule
Job ID & Script & Payload & Dependency \\
\midrule
27711 & \texttt{build\_container.sbatch} & build + verify \texttt{cacose.sif} & --- \\
27712 & \texttt{train\_benchmark\_array.sbatch} & \texttt{configs/cora.yaml}, seeds 0--9 & \texttt{afterok:27711} \\
27713 & \texttt{train\_benchmark\_array.sbatch} & \texttt{configs/mutag.yaml}, seeds 0--9 & \texttt{afterany:27712,afterok:27711} \\
27714 & \texttt{train\_benchmark\_array.sbatch} & \texttt{configs/chameleon.yaml}, seeds 0--9 & \texttt{afterany:27713,afterok:27711} \\
\bottomrule
\end{tabular}
\caption{The four jobs, submitted by one invocation of \texttt{scripts/submit\_all\_benchmarks.sh}.
Each array carries \texttt{--array=0-9\%3}, so at most three seeds run concurrently; the
\texttt{afterany} chain between datasets keeps that cap global rather than per-array.}
\end{table}

\section{The failed run that preceded this one}

The first attempt at this batch produced no data at all, and the reason is worth recording because
it was silent.

\texttt{APPTAINER\_TMPDIR} had been pointed at node-local \texttt{/tmp} to suppress a cosmetic
\texttt{nodev} warning about \texttt{/data}. On these nodes \texttt{/tmp} is \texttt{tmpfs}, i.e.\
RAM. \texttt{mksquashfs} unpacks the entire uncompressed image into that scratch directory before
compressing it, so the unpack was charged against the build job's memory cgroup. It ran out, and
\emph{the build did not fail}: it wrote a \texttt{.sif} with a valid header and a truncated
squashfs payload. \texttt{apptainer inspect} succeeded on that image, so the usual sanity check
reported it healthy, while every \texttt{apptainer exec} --- including \texttt{exec /bin/echo} ---
died with \texttt{input/output error}. All 30 benchmark tasks failed identically.

Two things made the failure durable rather than merely annoying:

\begin{enumerate}
\item The definition's \texttt{\%test} section passed. \texttt{\%test} runs against the build tree,
      not the finished \texttt{.sif}, so it cannot detect a truncated image and its passing was
      false assurance.
\item The build wrote its \texttt{.def.sha256} sidecar \emph{before} verifying the image. The hash
      matched the definition, so the freshness check in the submission path reported the container
      up to date and would have reused the broken image indefinitely.
\end{enumerate}

Both are fixed in the commit recorded above: scratch moved to the NAS, the sidecar is now written
only after the finished image has been executed successfully, and a verification failure exits
non-zero so that \texttt{--kill-on-invalid-dep} cancels the dependent sweeps. The build log for job
27711 ends with \texttt{Verified: yes}, which under the new ordering is a statement about the
finished artefact rather than about the build tree.

\section{Protocol}

One command was issued:

\begin{verbatim}
scripts/submit_all_benchmarks.sh
\end{verbatim}

which checked the container once, submitted the build, and chained the three sweeps behind it. No
task was resubmitted, no seed was rerun, and no configuration was edited while the batch was in
flight. That last point is stated explicitly because an earlier MUTAG measurement was contaminated
exactly that way (see \texttt{REPRO\_REPORT.md} \S5.5), and the discipline is now part of the
protocol rather than an assumption.

Each task runs \texttt{scripts.run\_single\_experiment} inside the container under
\texttt{--containall}, with only three bind mounts: the repository, the user's own \texttt{/data}
tree, and node-local \texttt{/tmp}. Datasets were prefetched in a prior job, because compute nodes
have no outbound network and a first-use download would fail after scheduling.

\subsection{Configurations}

Resolved configurations, as recorded inside each result JSON. Config hashes exclude the seed, so a
hash identifies the experiment and the ten seeds are replicates of it.

\begin{table}[!htbp]
\centering
\begin{tabular}{lrrr}
\toprule
 & Cora & Chameleon & MUTAG \\
\midrule
Config hash & \texttt{4c87d394} & \texttt{78dd98a6} & \texttt{ebbbb2c9} \\
Task & node clf. & node clf. & graph clf. \\
Provider & Planetoid & WikipediaNetwork & TUDataset \\
Split & 1208/500/500 fixed & 0.48/0.32/0.20 strat. & 0.80/0.10/0.10 strat. \\
Split sizes (nodes/graphs) & 1208/500/500 & 1093/729/455 & 150/18/20 \\
$\delta$ / \caef{} mode & 3 / single & 3 / single & 3 / single \\
Backbone / pooling & gcn / sagpool & gcn / sagpool & gcn / sagpool \\
Readout & mean$\|$max & mean$\|$max & \textbf{sum} \\
Attention heads & 2 & 2 & 1 \\
Dropout & 0.5 & 0.5 & 0.0 \\
Epochs / patience & 250 / 50 & 250 / 50 & 100 / 25 \\
\bottomrule
\end{tabular}
\caption{Learning rate $2.5\times10^{-3}$, weight decay $10^{-4}$, hidden width 128, $d_S=128$,
pooling ratio 0.5, batch size 32 throughout. Every value not stated by the paper is logged as an
ambiguity in the specification's \S3.}
\end{table}

\section{Results}

\subsection{Headline}

\begin{table}[!htbp]
\centering
\begin{tabular}{lrrrrl}
\toprule
Dataset & Paper & Accept at & This run & $\Delta$ & Outcome \\
\midrule
Cora      & 85.00 & $\geq 83.5$ & $84.64 \pm 1.75$  & $-0.36$ & PASS \\
Chameleon & 68.99 & $\geq 66.5$ & $67.19 \pm 2.74$  & $-1.80$ & PASS \\
MUTAG     & 76.99 & $\geq 74.0$ & $80.50 \pm 11.65$ & $+3.51$ & PASS \\
\bottomrule
\end{tabular}
\caption{Test accuracy, mean $\pm$ sample standard deviation over seeds 0--9. Acceptance
thresholds were fixed in the specification before any run and were not revised afterwards.}
\end{table}

\subsection{Per-seed results}

The tables below are the complete record: nothing is excluded, and no seed was discarded. Macro-F1
is reported alongside accuracy because on MUTAG the two diverge in a way that accuracy alone
conceals.

\begin{table}[!htbp]
\centering
\begin{tabular}{rrrrrrr}
\toprule
\multicolumn{7}{c}{\textbf{Cora} (job 27712, hash \texttt{4c87d394})} \\
\midrule
Seed & Test acc & Macro-F1 & Best val & Epochs & Best epoch & Time (s) \\
\midrule
0 & 85.00 & 82.49 & 84.20 & 70 & 20 & 9.7 \\
1 & 85.60 & 83.75 & 88.00 & 71 & 21 & 9.9 \\
2 & 82.00 & 79.16 & 86.60 & 71 & 21 & 9.7 \\
3 & 83.80 & 83.18 & 85.40 & 131 & 81 & 15.6 \\
4 & 84.00 & 82.57 & 86.20 & 86 & 36 & 11.2 \\
5 & 87.40 & 86.42 & 86.00 & 77 & 27 & 10.2 \\
6 & 83.80 & 82.43 & 86.20 & 72 & 22 & 9.8 \\
7 & 85.80 & 83.59 & 83.00 & 74 & 24 & 10.2 \\
8 & 82.40 & 81.95 & 85.20 & 75 & 25 & 9.3 \\
9 & 86.60 & 84.89 & 86.60 & 75 & 25 & 9.7 \\
\midrule
Mean & $84.64 \pm 1.75$ & $83.04 \pm 1.91$ & 85.74 & 80.2 & 30.2 & 10.5 \\
\bottomrule
\end{tabular}
\end{table}

\begin{table}[!htbp]
\centering
\begin{tabular}{rrrrrrr}
\toprule
\multicolumn{7}{c}{\textbf{Chameleon} (job 27714, hash \texttt{78dd98a6})} \\
\midrule
Seed & Test acc & Macro-F1 & Best val & Epochs & Best epoch & Time (s) \\
\midrule
0 & 64.18 & 64.13 & 67.76 & 105 & 55 & 108.4 \\
1 & 67.69 & 67.76 & 68.18 & 81 & 31 & 86.9 \\
2 & 69.89 & 69.83 & 69.55 & 162 & 112 & 164.6 \\
3 & 67.25 & 67.43 & 69.00 & 104 & 54 & 109.6 \\
4 & 61.32 & 61.25 & 68.72 & 85 & 35 & 89.6 \\
5 & 68.79 & 68.65 & 68.45 & 124 & 74 & 128.3 \\
6 & 67.91 & 67.64 & 69.41 & 99 & 49 & 105.2 \\
7 & 67.47 & 67.18 & 69.68 & 135 & 85 & 138.0 \\
8 & 66.59 & 66.43 & 67.22 & 134 & 84 & 133.5 \\
9 & 70.77 & 70.60 & 70.23 & 187 & 137 & 194.0 \\
\midrule
Mean & $67.19 \pm 2.74$ & $67.09 \pm 2.72$ & 68.82 & 121.6 & 71.6 & 125.8 \\
\bottomrule
\end{tabular}
\end{table}

\begin{table}[!htbp]
\centering
\begin{tabular}{rrrrrrr}
\toprule
\multicolumn{7}{c}{\textbf{MUTAG} (job 27713, hash \texttt{ebbbb2c9})} \\
\midrule
Seed & Test acc & Macro-F1 & Best val & Epochs & Best epoch & Time (s) \\
\midrule
0 & 85.00 & 82.91 & 77.78 & 46 & 21 & 13.3 \\
1 & 55.00 & 54.89 & 83.33 & 33 & 8 & 10.5 \\
2 & 95.00 & 94.67 & 88.89 & 52 & 27 & 14.7 \\
3 & 85.00 & 84.00 & 83.33 & 32 & 7 & 10.2 \\
4 & 75.00 & 73.33 & 83.33 & 35 & 10 & 10.6 \\
5 & 95.00 & 94.67 & 83.33 & 33 & 8 & 9.9 \\
6 & 85.00 & 81.19 & 77.78 & 27 & 2 & 9.3 \\
7 & 75.00 & 74.94 & 83.33 & 33 & 8 & 10.5 \\
8 & 80.00 & 78.02 & 100.00 & 78 & 53 & 20.1 \\
9 & 75.00 & 68.65 & 94.44 & 29 & 4 & 8.9 \\
\midrule
Mean & $80.50 \pm 11.65$ & $78.73 \pm 11.89$ & 85.56 & 39.8 & 14.8 & 11.8 \\
\bottomrule
\end{tabular}
\end{table}

Every run stopped early; none exhausted its epoch budget. On Cora and Chameleon this is
unremarkable. On MUTAG it is a caveat, discussed in \S\ref{sec:threats}.

\subsection{Decomposition statistics}

These are structural facts about the decomposition, independent of training, and they determine
both the cost and --- on MUTAG --- whether the method's own contribution is exercised at all.

\begin{table}[!htbp]
\centering
\begin{tabular}{lrrrr}
\toprule
Dataset & $k_{\max}$ & Subgraphs & Isolated nodes & Parameters \\
\midrule
Cora      & 4  & 4              & 0 & 1{,}032{,}207 \\
Chameleon & 63 & 50             & 0 & 17{,}462{,}889 \\
MUTAG     & 2  & 376 (2/graph)  & 0 & 118{,}150 \\
\bottomrule
\end{tabular}
\caption{Chameleon's 63 core levels yield roughly 50 separate GCN and pooling branches and 17.5M
parameters fitted to 2277 nodes. This was flagged as an over-parameterisation risk in the
specification before any run; it reproduced regardless, and \texttt{share\_weights = true} was
never required.}
\end{table}

MUTAG's $k_{\max} = 2$ is below $\delta = 3$, so \caef{} never fires on any molecule in the
dataset: the decomposition reduces to plain $k$-core and cross-subgraph attention operates over two
vectors. MUTAG therefore contributes almost nothing as evidence for the paper's central claim,
whatever its accuracy.

\subsection{Cost}

\begin{table}[!htbp]
\centering
\begin{tabular}{lrrr}
\toprule
Dataset & s / seed & Total (10 seeds) & Peak parameters \\
\midrule
Cora      & 10.5  & 1.8 min  & 1.0M \\
MUTAG     & 11.8  & 2.0 min  & 0.1M \\
Chameleon & 125.8 & 21.0 min & 17.5M \\
\midrule
All three & --- & 24.7 min compute, 11.8 min wall & --- \\
\bottomrule
\end{tabular}
\caption{Compute time is the sum over tasks; wall-clock is shorter because three seeds run
concurrently. The whole reproduction costs under half an hour of one GTX 1080 Ti.}
\end{table}

\section{Reproducibility of this record}
\label{sec:determinism}

The runner seeds Python, NumPy, torch, CUDA and PyG, and enables
\texttt{torch.use\_deterministic\_algorithms(True, warn\_only=True)}. Despite that, \textbf{these
GPU numbers are not bit-reproducible}, and the logs say so explicitly: every task emitted

\begin{quote}\small
\texttt{UserWarning: Deterministic behavior was enabled ... but this operation is not
deterministic because it uses CuBLAS and you have CUDA >= 10.2. To enable deterministic behavior
... set CUBLAS\_WORKSPACE\_CONFIG=:4096:8}
\end{quote}

raised from \texttt{GCNConv}'s linear layer, the multi-head attention \texttt{bmm} calls, and the
backward pass. The environment variable was not set for this run, so a rerun at the same commit may
differ in the low-order digits. The fix is a single \texttt{--env} in the array script and is
tracked as follow-up work; it was deliberately not applied retroactively, because doing so would
make this document describe a run that did not happen.

The \texttt{warn\_only=True} setting is itself a considered choice: several PyG kernels have no
deterministic implementation, and hard-failing would block the run rather than make it
reproducible.

The practical size of that nondeterminism can be bounded from the record. An earlier cluster GPU
run at commit \texttt{1aecbcfdaec4}, under the \emph{same} config hashes \texttt{4c87d394} and
\texttt{78dd98a6}, gave Cora $84.62 \pm 1.91$ and Chameleon $67.25 \pm 2.24$, against $84.64$ and
$67.19$ here: a drift of $0.02$ and $0.06$ points. The configuration is identical and the
intervening commits touched the submission scripts rather than the model or training loop, so that
drift is an upper bound on the nondeterminism --- an upper bound rather than a measurement of it,
because that earlier run's node and GPU model were not recorded, and hardware may therefore be
contributing too. Either way it is an order of magnitude below the seed standard deviations and
affects no conclusion drawn from these means.

That the earlier run's hardware went unrecorded is itself a defect in the previous bookkeeping, and
is why the provenance table in \S1 exists.

This is deliberately a weaker claim than the one an earlier draft of \texttt{REPRO\_REPORT.md}
made. That draft asserted CPU and GPU agreed \emph{exactly}, which was both false and, as stated,
unverifiable from the artefacts retained: no CPU run of these two configurations survives in the
result store, so no cross-platform claim is made here at all.

\section{Threats to validity}
\label{sec:threats}

Stated here rather than left for a reader to notice.

\paragraph{MUTAG's test set has 20 graphs.} One graph is 5 percentage points. Every per-seed
accuracy in the MUTAG table is a multiple of 5, the seed spread runs from 55.0 to 95.0, and the
standard deviation of 11.65 is a property of the split size, not of the method. The same arithmetic
is what shows the paper's reported 76.99 to be unreachable under its own stated 80/10/10 protocol,
since a 10-seed mean over $n$ test graphs must be a multiple of $1/(10n)$.

\paragraph{MUTAG's validation set has 18 graphs, and early stopping runs on it.} Best-epoch values
of 2, 4, 7 and 8 appear in the table: on several seeds the validation score peaked almost
immediately and patience 25 then halted training around epoch 30 of a 100-epoch budget. Model
selection on 18 graphs is close to arbitrary, and this, rather than optimisation, plausibly drives
much of the MUTAG variance.

\paragraph{The MUTAG readout is our choice, not the paper's.} The paper states only ``READOUT''.
\texttt{sum} was selected by measurement --- it is worth roughly 13 accuracy points and 18 F1
points over mean$\|$max --- and the selection is defensible (sum is the only injective multiset
readout, and molecular labels depend on substructure counts), but it remains a choice the paper
does not license. Cora and Chameleon use the conventional mean$\|$max and required no such
decision.

\paragraph{Single node, single GPU model.} All 30 tasks ran on \texttt{pleiades-1-3} on one GTX
1080 Ti. Nothing here establishes behaviour on other hardware.

\paragraph{Ten seeds is few.} With Chameleon's spread of 61.3--70.8, the standard error on the mean
is roughly 0.9 points, so the $-1.80$ gap to the paper is around two standard errors: suggestive of
a real shortfall, not conclusive.

\section{Log inventory}

\texttt{/results/} is excluded from version control, so the tables above are the archival form of
this run. The raw artefacts, should they still exist, are:

\begin{table}[!htbp]
\centering
\begin{tabular}{ll}
\toprule
Artefact & Path \\
\midrule
Result JSON (30) & \texttt{results/<dataset>/<config\_hash>/seed<NN>.json} \\
Task stdout (30) & \texttt{logs/cacose\_<array\_job>\_<task>.out} \\
Task stderr (30) & \texttt{logs/cacose\_<array\_job>\_<task>.err} \\
Build log & \texttt{logs/cacose-build\_27711.\{out,err\}} \\
\bottomrule
\end{tabular}
\caption{Relative to \texttt{/data/\$USER/CaCoSE/} on the cluster. Each result JSON carries the
commit, the fully resolved configuration, the split sizes, the torch and PyG versions, and the
device, so any row in this document can be traced to its origin without the surrounding logs.}
\end{table}

Every task's stderr contains only the CuBLAS determinism warnings of
\S\ref{sec:determinism} plus a note that \texttt{torch-scatter} is absent (an optional accelerator,
not a correctness dependency). The files are byte-identical in size within a dataset, which is what
first indicated that no task had failed.

\section{Regenerating this run}

\begin{verbatim}
git clone https://github.com/M4-on-Github/GNN_CaCoSE.git
cd GNN_CaCoSE && git checkout 2147b500cd12
sbatch slurm/download_datasets.sbatch     # once; compute nodes have no network
scripts/submit_all_benchmarks.sh          # builds the container if needed, then all three
CACOSE_OUT=/data/$USER/CaCoSE python3 -m scripts.aggregate_benchmark_results
\end{verbatim}

The aggregation script needs neither torch nor a GPU and reproduces the headline table from the
JSON files alone. Subject to \S\ref{sec:determinism}, expect agreement to a few hundredths rather
than exact equality.

\end{document}
